

\documentclass{svmult}



\newcommand{\opteqn}[1]{$#1$}
\newcommand{\optpar}{}  % change to null and linefeed
\newcommand{\mysubsection}[1]{\smallskip\noindent{\bf \emph{#1}.}}


\usepackage{makeidx}
\usepackage{graphicx}
\usepackage{amssymb}

\newcommand{\eq}[2]{\begin{equation}\label{#1}#2\end{equation}}

\newcommand{\pic}[2]{\includegraphics[scale=#1]{\FIGDIR/#2}}
\newcommand{\picc}[2]{\begin{center}\pic{#1}{#2}\end{center}}

\newcommand{\comment}[1]{}
\newcommand{\edo}{\end{document}}


\newcommand{\R}{{\mathbb R}}  %ams bold

\newcommand{\abs}[1]{\left\vert #1 \right\vert}
\newcommand{\norm}[1]{\left\Vert #1 \right\Vert}

\newcommand{\be}[1]{\begin{equation}\label{#1}}
\newcommand{\ee}{\end{equation}}

\newcommand{\halmos}{\rule{1ex}{1.4ex}}

\newcommand{\mybox}{\hfill $\Box$} %put \qed at right margin (white square)

\newcommand{\beq}{\begin{eqnarray}}
\newcommand{\eeq}{\end{eqnarray}}
\newcommand{\beqn}{\begin{eqnarray*}}
\newcommand{\eeqn}{\end{eqnarray*}}
\newcommand{\bi}{\begin{itemize}}
\newcommand{\ei}{\end{itemize}}
\newcommand{\ben}{\begin{enumerate}}
\newcommand{\een}{\end{enumerate}}

\newcommand{\st}{\, | \,}
\newcommand{\inter}{{\rm int}}
\newcommand{\co}{{\rm co}}
\newcommand{\col}{{\rm col}}

\renewcommand{\thefootnote}{\fnsymbol{footnote}}







\newcommand{\cc}{c}
\newcommand{\kk}{k}
\newcommand{\cckT}{ckT}
\newcommand{\pp}{p}
\newcommand{\norma}{1}
\newcommand{\normb}{2}

\newcommand{\perorb}{\hat{x}}
\newcommand{\mypmatrix}[1]{\left(\begin{array}{cccccccccccc}#1\end{array}\right)}

\newcommand{\absv}[1]{\left\vert #1 \right\vert}

\title*{Contractive Systems with Inputs}

\titlerunning{Contractive Systems with Inputs}

\author{Eduardo D. Sontag}
\institute{Department of Mathematics\\
Rutgers University\\
USA\\
\texttt{sontag@math.rutgers.edu}}

\begin{document}

\maketitle

\medskip

\centerline{$\framebox{\bf Dedicated to Y. Yamamoto on the occasion of his 60th birthday}$}

\bigskip

\subsection*{Abstract}

Contraction theory provides an elegant way of analyzing the behaviors of
systems subject to external inputs.  Under sometimes easy to check hypotheses,
systems can be shown to have the incremental stability property that all
trajectories converge to a unique solution.  This property is especially
interesting when forcing functions are periodic (a globally attracting limit
cycle results), as well as in the context of establishing synchronization
results.  The present paper provides a self-contained introduction to some
basic results, with a focus on contractions with respect to non-Euclidean
metrics.

\section{Introduction}


The most common approach to analyzing global stability properties of nonlinear
dynamical systems is through Lyapunov functions.
However, in many applications, Lyapunov functions are not always
easy to find, especially if steady states are not known \emph{a priori}.
Remarkably, a stronger property than stability, namely the \emph{contraction}
(or incremental stability) requirement that all solutions should converge
(exponentially) towards each other, is sometimes easier to work with.
Contractive dynamics result when the logarithmic norm, or matrix measure, of
the Jacobian of the vector field is uniformly negative on the state space.
Different norms are appropriate to different problems, just as different
Lyapunov functions have to be carefully picked.
Non-Euclidean norms have been found to be useful in the study of many
bio-molecular problems, see for example~\cite{russo_dibernardo_sontag}.

The study of contractions in the context of stability theory dates back at
least to the work of Demidovich (\cite{demidovich}), who established the basic
convergence results with respect to Euclidean norms, and independently to
Yoshizawa (\cite{yoshizawa1,yoshizawa2}); see~\cite{pavlov} for a historical
discussion. 
In control theory, contraction theory has been popularized and extended by 
Slotine and coworkers, see for instance 
\cite{Loh_Slo_00}, \cite{Jou_04_a}, \cite{Wan_Slo_05}
where applications to nonlinear control, observer problems, and synchronization
and consensus problems in complex networks have been developed,
as well as by Nijmejer and coworkers in the context of nonlinear regulator
problems, see for example \cite{pavlov_book}.
In this latter work, the authors use the phrase ``convergent dynamics''
to refer to property that there exists a (necessarily unique) globally
asymptotically stable solution to which all other solutions converge.

This paper gives a self-contained exposition, with simple proofs, of some
basic results in contraction theory, It seems difficult to find such simple
proofs in the literature, particularly for contractions with respect to
non-Euclidean norms.  We emphasize that the presentation is expository, and no
substantial new results on contraction theory are claimed.

Definitions and statements of the main results are provided in
Section~\ref{sec:defs}, and proofs are given in Section~\ref{sec:proofs}.

Section~\ref{sec:synch} briefly discusses the application of contraction
theory to the synchronization of coupled identical dynamical systems,
following an idea of Slotine and collaborators (``virtual systems'').  Also
discussed there is a minor extension in which simultaneous convergence, not
merely synchronization, is achieved.

For periodically forced contractive systems, globally attracting limit cycles
arise, a sort of ``entrainment'' property.
Such a property is false for general systems that have a well-defined steady-state
response to constant inputs, for which even chaotic behavior may arise under
periodic forcing (\cite{arxiv_chaos_2009}).

In closing this introduction, we remark that a modern approach to contractive
dynamics steps away from the consideration of Jacobians, and defines
contraction properties by means of ``logarithmic Lipschitz constants''
directly associated to the vector field.  This elegant approach, nicely
surveyed in~\cite{soderlind}, is powerful and intuitive, and allows immediate
generalizations to infinite-dimensional problems.  However, in order to verify
the property for particular examples, Jacobians must still be employed.


\section{Definitions and Statements of Main Results}
\label{sec:defs}

We consider in this paper systems of ordinary differential equations,
generally time-dependent:
\be{eqn:gensys}
\dot x = f(t,x)
\ee
defined for $t\in [0,\infty )$ and $x\in C$, where $C$ is a subset of $\R^n$.
It will be assumed that $f(t,x)$ is differentiable on $x$, and that
$f(t,x)$, as well as the Jacobian of $f$ with respect to $x$,
denoted as $J(t,x) = \frac{\partial f}{\partial x}(t,x)$, are continuous in $(t,x)$. 
In applications of the theory, it is often the case that $C$ will
be a closed set, for example given by non-negativity constraints on variables
as well as linear equalities representing mass-conservation laws.
For a non-open set $C$, differentiability in $x$ means that the vector field
$f(t,\cdot )$ can be extended as a differentiable function to some open set which
includes $C$, and the continuity hypotheses with respect to $(t,x)$ hold on
this open set.

We denote by
$
\varphi(t,s,\xi )
$
the value of the solution $x(t)$ at time $t$ of the differential
equation~(\ref{eqn:gensys}) with initial value $x(s)=\xi $.
It is implicit in the notation that $\varphi(t,s,\xi )\in C$ (``forward invariance''
of the state set $C$).
This solution is in principle defined only on some interval $s\leq t<s+\varepsilon $,
but we will assume that $\varphi(t,s,\xi )$ is defined for all $t\geq s$.
Conditions which guarantee such a ``forward-completeness'' property
are often satisfied in applications, for example whenever
the set $C$ is closed and bounded, or whenever the vector field $f$ is
bounded.  (See for example Appendix C in~\cite{mct} for more discussion, as well
as~\cite{angeli-sontag-fc} for a characterization of the forward completeness
property.) 
Under the stated assumptions, the function $\varphi$ is jointly differentiable
in all its arguments (this is a standard fact on well-posedness of
differential equations, see for example Appendix C in~\cite{mct}).

We recall (see for instance~\cite{michelbook} or
\cite{desoer-vidyasagar-siam}) that, given a vector norm on Euclidean space
($\abs{\cdot }$), with its induced matrix norm $\norm{A}$, the associated
\emph{matrix measure} $\mu $ is defined as the directional derivative of the
matrix norm in the direction of $A$ and evaluated at the identity matrix, that
is:
\opteqn{
\mu (A) \,:=\;
\lim_{h \searrow 0} \frac{1}{h} \left(\norm{I+hA}-1\right).
}
For example, if $\abs{\cdot }$ is the standard Euclidean 2-norm, then
$\mu (A)$ is the maximum eigenvalue of the symmetric part of $A$.
Matrix measures, also known as ``\emph{logarithmic norms}'', were
independently introduced by Germund Dahlquist and Sergei Lozinskii in 1959,
\cite{dahlquist,lozinskii}.
The limit is known to exist, and the convergence is monotonic,
see~\cite{strom,dahlquist}.

\begin{definition}
The system (\ref{eqn:gensys}), or the time-dependent vector field $f$, is
said to be
\emph{infinitesimally contracting} on a set $C \subseteq \R ^n$ if 
there exists some norm in $C$, with associated matrix measure $\mu $, such that,
for some constant $c>0$ (the \emph{contraction rate}), it holds that:
\begin{equation} \label{eqn:contrcond}
\mu \left(J\left(x,t\right)\right) \le - \cc, \quad \forall x \in C, \quad \forall t \ge 0.
\end{equation}
\end{definition}

The key result is that infinitesimal contractivity implies global contractivity:

\begin{theorem}\label{theo:main}
Suppose that $C$ is a convex subset of $\R^n$ and that $f(t,x)$ is
infinitesimally contracting with contraction rate $\cc$.
Then, for every two solutions $x(t)=\varphi(t,0,\xi )$ and $z(t)=\varphi(t,0,\zeta )$
of~(\ref{eqn:gensys}), it holds that:
\be{eqn:contract}
\abs{x(t)-z(t)} \;\leq \; e^{-\cc t} \abs{\xi -\zeta }, 
\quad\quad\forall\,t\ge 0\,.
\ee
\end{theorem}

If ${\cal A}$ is a non-empty forward-invariant set for the dynamics,
then every solution must approach ${\cal A}$.  Indeed, take any $\zeta \in {\cal A}$ and any
trajectory $x(t)=\varphi(t,0,\xi )$; then, with $z(t)=\varphi(t,0,\zeta )$,
$\mbox{dist}\,(x(t),{\cal A})\leq \abs{x(t)-z(t)} \;\leq \; e^{-\cc t} \abs{\xi -\zeta }\rightarrow 0$
as $t\rightarrow \infty $.  In particular, if an equilibrium exists, then it must be
unique and globally asymptotically stable, and the same is true for periodic
orbits.  More interestingly, periodic orbits are assured to exist if
the vector field is periodic, as would happen for a system with inputs
$\dot x=f(x,u)$ under a periodic input $u(\cdot )$.  We discuss this next.

Given a number $T>0$, we will say that system~(\ref{eqn:gensys}) is
\emph{$T$-periodic} if it holds that 
\opteqn{
f(t+T,x) \,=\, f(t,x) 
\quad\forall\,t\geq 0,\,x\in C\,.
}
Notice that a system $\dot x=f\left(x,u(t)\right)$ with input $u(t)$ is
$T$-periodic $u(t)$ is itself a periodic function of period $T$. 
\optpar
The basic theoretical result about periodic orbits is as follows.

\begin{theorem}\label{theorem:contraction}
Suppose that:
\begin{itemize}
\item $C$ is a closed convex subset of $\R^n$;
\item $f$ is infinitesimally contracting with contraction rate $\cc$;
\item $f$ is $T$-periodic.
\end{itemize}
Then, there is a unique periodic solution $\perorb(t):[0,\infty )\rightarrow C$ of~(\ref{eqn:gensys})
of period $T$ and,
for every solution $x(t)$, it holds that $\abs{x(t)-\perorb(t)}\rightarrow 0$ as $t\rightarrow \infty $.
\end{theorem}


Cascades of contractive systems are again contracting.  To state this fact
precisely, let us consider a system of the following form:
\beqn
\dot  x &=& f(t,x)\\
\dot  y &=& g(t,x,y)
\eeqn
where $x(t)\in C_1\subseteq \R^{n_1}$ and $y(t)\in C_2\subseteq \R^{n_2}$ for
all $t$.
We write the Jacobian of $f$ with respect to $x$ as
$A(t,x) = \frac{\partial f}{\partial x}(t,x)$,
the Jacobian of $g$ with respect to $x$ as
$B(t,x,y) = \frac{\partial g}{\partial x}(t,x,y)$,
and the Jacobian of $g$ with respect to $y$ as
$C(t,x,y) = \frac{\partial g}{\partial y}(t,x,y)$,

When we say that $\dot y=g(t,x,y)$ is \emph{infinitesimally contracting
when $x$ is viewed as a parameter} we mean that, with respect
to some norm $\abs{\cdot }$), there is an estimate
$\mu (C(t,x,y))\leq  -\cc_2<0$
for all $x\in C_1$, $y\in C_2$ and all $t\geq 0$.

\begin{theorem}\label{theo:cascades}
Suppose that:
\bi
\item
the system $\dot x=f(t,x)$ is infinitesimally contracting;
\item
the system $\dot y=g(t,x,y)$ is infinitesimally contracting when $x$ is
viewed as a parameter;
\item
the mixed Jacobian $B(t,x,y)$ is bounded.
\ei
Then, the cascaded system is infinitesimally contracting.
\end{theorem}

The basic contraction property insures that any solutions of $\dot x=f(t,x)$
exponentially converge to each other.  The following result provides
a ``robustness margin'' that says that any solution of the original system
and any solution of a perturbed system $\dot x=f(t,x)+h(t)$ also
exponentially converge to each other, provided that $h(t)\rightarrow 0$ exponentially.
This is a ``converging-input converging output'' property that provides a weak
type of input-to-state stability.

\begin{theorem}
\label{theo:fhc}
Assume that the system $\dot x=f(t,x)$ is infinitesimally contracting.
Let $h(t)$ be a vector function satisfying
\opteqn{
\abs{h(t)}\leq Le^{-kt} \quad \forall\,t\geq 0
}
for some $k>0$ and $L\geq 0$,
Then, there exist constants $\ell>0$ and $\kappa $ such that the following property
holds:
For any solution $x(t)=\varphi(t,0,\xi )$ of the system $\dot x=f(t,x)$,
and any solution $z(t)=\varphi(t,0,\zeta )$ of the system $\dot x=f(t,x)+h(t)$,
\be{eqn:estimatefh}
\abs{x(t)-z(t)} \leq  e^{-\ell t} \left(\kappa  +  \abs{\xi -\zeta }\right)
\ee
for all $t\geq 0$.
\end{theorem}

In general, the constant $\kappa $ cannot be dropped from the estimate
in Theorem~\ref{theo:fhc}.
Indeed, consider this counterexample: compare the solutions $x(t)=0$ and
$z(t)=te^{-t}$ of $\dot x=-x$ and $\dot x=-x+e^{-t}$ with $\xi =\zeta =0$ respectively.

Observe that any solutions of $\dot x=f(t,x)+h_1(t)$ and $\dot x=g(t,x)+h_2(t)$
will also converge to each other, if $h_1$ and $h_2$ satisfy the properties
for $h$ in Theorem~\ref{theo:fhc}, since they both converge to any solution
of the system with no $h$.

\section{Proofs of Main Results}
\label{sec:proofs}

\mysubsection{Proof of Theorem~\ref{theo:main}}
We give the proof in a generalized form, in which convexity is replaced by a
weaker constraint on the geometry of the space, but the estimate on
trajectories is potentially weaker than in the convex case.

Let $K>0$ be any positive real number and assume that a norm in $\R^n$ has
been chosen.
We will say that a subset $C\subset\R^n$ is \emph{$K$-reachable} if, for any
two points $x_0$ and $y_0$ in $C$ there is some continuously differentiable
curve $\gamma : \left[  0, 1 \right] \rightarrow C$ 
such that:
$\gamma \left(0\right) = x_0$;
$\gamma \left(1\right)=y_0$;
$\abs{\gamma '  \left(r\right)} \le K \abs{y_0 - x_0}$ for all $r\in [0,1]$. 
For convex sets $C$, we may pick $\gamma (r)=x_0+r(y_0-x_0)$, so 
$\gamma '(r)= y_0-x_0$ and we can take $K=1$.  Thus, convex sets are
$1$-reachable, and it is easy to show that the converse holds as well.

Note that a set $C$ is $K$-reachable for some $K$ if and only if the length
of a minimal-length (geodesic) smooth path connecting any two points $x$
and $y$ in $C$ and parametrized by arc length, is bounded by some multiple
$K_0$ of the Euclidean norm $\abs{y-x}_2$. Indeed, re-parametrizing to a path
$\gamma $ defined on $\left[0,1\right]$, we have:
\opteqn{
\abs{\gamma ' \left(r\right)}_2 \le K_0 \abs{y-x}_2.
}
Since in finite dimensional spaces all norms are equivalent, a
suitable $K$ as in the above estimate exists.

\begin{lemma}\label{theo:k-reachable}
Suppose that $C$ is a $K$-reachable subset of $\R^n$ and that $f(t,x)$ is
infinitesimally contracting with contraction rate $\cc$.
Then, for every two solutions $x(t)=\varphi(t,0,\xi )$ and $z(t)=\varphi(t,0,\zeta )$
it holds that:
\be{eqn:weak}
\abs{x(t)-z(t)} \leq  K e^{-\cc t} \abs{\xi -\zeta } 
\quad\quad\forall\,t\ge 0\,.
\ee
\end{lemma}

Observe that Theorem~\ref{theo:main}
follows trivially from Lemma \ref{theo:k-reachable}, since for convex sets
we may pick $K=1$.

\begin{proof}
Given any two points $x\left(0\right)= \xi $ and $z\left(0\right)= \zeta $ in
$C$, pick a smooth curve $\gamma : \left[ 0, 1\right] \rightarrow C$, such
that $\gamma \left(0\right)=\xi $ and $\gamma \left(1\right)=\zeta $. Let  
$\psi \left(t,r\right)=\varphi(t,0,\gamma \left(r\right))$, that is,
the solution of system (\ref{eqn:gensys}) rooted at
$\psi \left(0,r\right) =\gamma \left(r\right)$, $r \in [0,1]$. 
Since $\varphi$ and $\gamma $ are continuously differentiable, also
$\psi \left(t,r\right)$ is continuously differentiable in both arguments.
We define
\opteqn{
w(t,r) := \frac{\partial \psi }{\partial r}(t,r).
}
It follows that
\[
\frac{\partial w}{\partial t}(t,r) = \frac{\partial}{\partial t}\left(\frac{\partial \psi }{\partial r}\right)=\frac{\partial}{\partial r}\left(\frac{\partial \psi }{\partial t}\right)=\frac{\partial}{\partial r}f(\psi \left(t,r\right),t).
\]
As
\opteqn{
\frac{\partial}{\partial r}f(\psi \left(t,r\right),t) = \frac{\partial f}{\partial x}(\psi \left(t,r\right),t)\frac{\partial \psi }{\partial r}(t,r),
}
\opteqn{
\frac{\partial w}{\partial t}(t,r) = J(\psi \left(t,r\right),t) w(t,r),
}
where $J(\psi \left(t,r\right),t) = \frac{\partial f}{\partial x}(\psi \left(t,r\right),t)$.
Appealing to  Coppel's inequality (see e.g.~\cite{Vid_93}), we have:
\begin{equation} \label{eq:coppelthm}
\abs{w(t,r)} \le \abs{w(0,r)}  e^{\int_{0}^t \mu \left( J\left(\tau \right)\right)d\tau }\leq K \abs{\xi -\zeta } e^{-\cc t},
\end{equation}
for all $x \in C$, $t\geq 0$, and $r \in [0,1]$. 
By the Fundamental Theorem of Calculus, we can write
\opteqn{
\psi \left(t,1\right)-\psi \left(t,0\right) = \int_0^1{w(t,s)} ds .
}
Hence, we obtain
\opteqn{
\abs{x(t)-z(t)} \leq \int_0^1{\abs{w(t,s)}}ds.
}
Now, using (\ref{eq:coppelthm}), the above inequality becomes:
\[
\abs{x(t)-z(t)} \le \int_0^1 \left(\abs{w(0,s)}  e^{\int_{0}^t \mu \left( J\left(\tau \right)\right)d\tau } \right) ds \le K \abs{\xi -\zeta }e^{-\cc t}.
\]
This completes the proof of the lemma.
\end{proof}

We remark that in some cases it might be possible to prove a strict contraction
($K=1$) even if the domain is not convex, by appealing to the deeper theory
of logarithmic Lipschitz constants (see~\cite{soderlind} for definitions and
details).  If the (lub) logarithmic Lipschitz constant $M[f]$ of the vector
field is $-c<0$, then an estimate~(\ref{eqn:contract}) holds.  
In general, $M[f]$ is an upper bound on the supremum of $\mu (J(t,x))$,
with equality to the supremum in the convex case.

\mysubsection{Proof of Theorem \ref{theorem:contraction}}
We assume now that the vector field $f$ is $T$-periodic.

\begin{remark}\label{rem:periodicf}
Periodicity implies that the initial
time is only relevant modulo $T$.  More precisely:
\be{eqn:periodicf}
\varphi(kT+t,kT,\xi ) = \varphi(t,0,\xi )
\ee
for all positive integers $k$, all $t\geq 0$, and all $x\in C$.
Indeed, let $z(s)=\varphi(s,kT,\xi )$, $s\geq kT$, and consider the function
$x(t)=z(kT+t)=\varphi(kT+t,kT,\xi )$, for $t\geq 0$.
So,
\[
\dot x(t) = \dot z(kT+t) = f(kT+t,z(kT+t)) = f(kT+t,x(t)) = f(t,x(t))\,,
\]
where the last equality follows by $T$-periodicity of $f$.
Since $x(0)=z(kT)=\varphi(kT,kT,\xi )=\xi $, it follows by uniqueness of
solutions that $x(t)=\varphi(t,0,\xi ) = \varphi \left(kT+t,kT,\xi \right)$, which is~(\ref{eqn:periodicf}).
As a corollary, we also have that
\be{eqn:periodicfc}
\varphi(kT+t,0,\xi ) = \varphi(kT+t,kT,\varphi(kT,0,\xi )) = \varphi(t,0,\varphi(kT,0,\xi ))
\ee
for all positive integers $k$, all $t\geq 0$, and all $x\in C$,
where the first equality follows from the semigroup property of solutions (see e.g. \cite{mct}), and
the second one from~(\ref{eqn:periodicf}) applied to
$\varphi(kT,0,\xi )$ instead of $\xi $.
\end{remark}

Define now
\opteqn{
P(\xi ) = \varphi(T,0,\xi ),
}
where $\xi = x \left(0\right) \in C$. 

\begin{lemma}\label{lemma:P}
$P^k(\xi ) = \varphi(kT,0,\xi )$ for all positive integers $k$ and $\xi \in C$.
\end{lemma}

\begin{proof}
We will prove the Lemma by recursion. In particular, the statement is true by definition when $k=1$.
Inductively, assuming it true for $k$, we have:
\[
P^{k+1}(\xi ) = 
P(P^k(\xi )) =
\varphi(T,0,P^k(\xi )) = 
\varphi(T,0,\varphi(kT,0,\xi )) =
\varphi(kT+T,0,\xi ).
\]
\end{proof}

\begin{theorem}\label{theo:contraction-general}
Suppose that:
\begin{itemize}
\item $C$ is a closed $K$-reachable subset of $\R^n$;
\item $f$ is infinitesimally contracting with contraction rate $\cc$;
\item $f$ is $T$-periodic;
\item $Ke^{-\cc T}<1$.
\end{itemize}
Then, there is an (unique) periodic solution $\perorb(t):[0,\infty )\rightarrow C$
of~(\ref{eqn:gensys}) having period $T$. Furthermore, every solution $x(t)$
converges to $\perorb(t)$, i.e. $\abs{x(t)-\perorb(t)}\rightarrow 0$ as $t\rightarrow \infty $. 
\end{theorem}

Theorem~\ref{theorem:contraction} is a corollary, because the assumption
$Ke^{-\cc T}<1$ in Theorem \ref{theo:contraction-general} is automatically
satisfied when the set $C$ is convex (i.e. $K=1$) and the system is
infinitesimally contracting.


\begin{proof}
Observe that $P$ is a contraction with factor $Ke^{-\cc T}<1$: 
$\abs{P(\xi )-P(\zeta )}\leq Ke^{-\cc T} \abs{\xi -\zeta }$
for all $\xi ,\zeta \in C$,
as a consequence of Theorem~\ref{theo:k-reachable}.
The set $C$ is a closed subset of $\R^n$ and hence complete as a metric space
with respect to the distance induced by the norm being considered.
Thus, by the contraction mapping theorem, there is a (unique) fixed
point $\bar \xi $ of $P$.
Let $\perorb(t):=\varphi(t,0,\bar \xi )$.
Since $\perorb(T)=P(\bar \xi )={\bar \xi }=\perorb(0)$, $\perorb(t)$ is a periodic orbit of
period $T$. 
Moreover, again by Theorem~\ref{theo:k-reachable},
we have that
$\abs{x(t)-\perorb(t)} \leq  K e^{-\cc t} \abs{\xi -\bar \xi }\rightarrow 0$.
Uniqueness is clear, since two different periodic orbits would be disjoint
compact subsets, and hence at positive distance from each other,
contradicting convergence. This completes the proof.
\end{proof}

Notice that, even in the non-convex case, the assumption $Ke^{-\cc T}<1$ can be
dropped, provided that we assert only the existence of (and global
convergence to) a unique periodic orbit, whose period is $kT$ for some integer
$k>1$.  Indeed, the vector field is also $kT$-periodic for any integer $k$.
Picking $k$ large enough so that $Ke^{-\cckT}<1$, we have the conclusion that
such an orbit exists, applying Theorem~\ref{theo:contraction-general}.

\mysubsection{Proof of Theorem~\protect{\ref{theo:cascades}}}
We assume that the system $\dot x=f(t,x)$ is infinitesimally contracting with
respect to a norm $\abs{\cdot }_\norma$, with contraction rate $\cc_1$, that is,
$\mu _\norma(A(t,x))\leq  -\cc_1$ for all $x\in C_1$ and all $t\geq 0$, where
$\mu _\norma$ is the matrix measure associated to $\abs{\cdot }_\norma$,
the system $\dot y=g(t,x,y)$ is infinitesimally contracting with respect to a
norm $\abs{\cdot }_{\normb}$ with contraction rate $\cc_2$, when $x$ is viewed as
a parameter in the second system, that is, $\mu _{\normb}(C(t,x,y))\leq  -\cc_2$
for all $x\in C_1$, $y\in C_2$ and all $t\geq 0$, where $\mu _{\normb}$ is
the matrix measure associated to $\abs{\cdot }_{\normb}$,
and that the mixed Jacobian $B(t,x,y)$ is bounded:
$\norm{B(t,x,y)}\leq \kk$, for all $x\in C_1$, $y\in C_2$ and all $t\geq 0$,
for some real number $k$, where ``$\norm{\cdot }$'' is the
operator norm induced by $\abs{\cdot }_{\norma}$ and $\abs{\cdot }_{\normb}$ on
linear operators $\R^{n_1} \rightarrow  \R^{n_2}$.

We need to show that, under these assumptions, the complete system is
infinitesimally contracting.  More precisely, pick any two positive numbers
$\rho _1$ and $\rho _2$ such that 
\opteqn{
\cc_1 \,-\, \frac{\rho _2}{\rho _1} \kk \;>\; 0
}
and let
\opteqn{
\cc := \min \left\{\cc_1 - \frac{\rho _2}{\rho _1} \kk, \cc_2\right\} \,.
}
We will show that $\mu (J)\leq  -\cc$, where $J$ is the full Jacobian:
\opteqn{\tiny 
J = 
\left[ 
\begin{array}{*{20}c}
A & 0 \\
B & C
\end{array} \right]\,,
}
with respect to the matrix measure $\mu $ induced by the following norm in
$\R^{n_1+n_2}$:
\opteqn{
\abs{(x_1,x_2)} = \rho _1 \abs{x_1}_{\norma} + \rho _2 \abs{x_2}_{\normb} \,.
}
Since
\opteqn{\tiny
(I+hJ)x = 
\left[ 
\begin{array}{*{20}c}
(I+hA)x_1  \\
hBx_1 + (I+hC)x_2
\end{array} \right]
}
for all $h$ and $x$, we have that, for all $h$ and $x$:
\beqn
\abs{(I+hJ)x}
&=&
\rho _1\abs{(I+hA)x_1} + \rho _2 \abs{hBx_1 + (I+hC)x_2}\\
&\leq &
\rho _1\abs{I+hA}\abs{x_1} +
\rho _2\abs{hB}\abs{x_1} +
\rho _2\abs{I+hC}\abs{x_2},
\eeqn
where from now on we drop subscripts for norms.
Pick now any $h>0$
and a unit vector $x$ (which depends on $h$)
such that $\norm{I+hJ} = \abs{(I+hJ)x}$.
Such a vector $x$ exists by the definition of induced matrix norm, and we note
that $1= \abs{x} = \rho _1\abs{x_1}_{\normb}+\rho _2\abs{x_2}_{\normb}$, by the definition of the norm
in the product space.
Therefore:
\beqn
\frac{1}{h} 
&&
\!\!\!\!\!
\left(\norm{I+hJ}-1\right)
\;=\;
\frac{1}{h} \left(\abs{(I+hJ)x}-\abs{x}\right)\\
&\leq &
\frac{1}{h} 
\left(
\rho _1\abs{I+hA}\abs{x_1} +
\rho _2\abs{hB}\abs{x_1} +
\rho _2\abs{I+hC}\abs{x_2}
- \rho _1 \abs{x_1} - \rho _2 \abs{x_2}
\right)\\
&=&
\frac{1}{h} 
\left(\abs{I+hA} -1 + \frac{\rho _2}{\rho _1}h\abs{B} \right) \rho _1\abs{x_1}
+
\frac{1}{h}
\left(\abs{I+hC} - 1 \right) \rho _2\abs{x_2}\\
&\leq &
\max\left\{
\frac{1}{h}\left(\abs{I+hA} -1\right) + \frac{\rho _2}{\rho _1}\kk
\,,\,
\frac{1}{h}\left(\abs{I+hC} - 1 \right)
\right\} \,,
\eeqn
where the last inequality is a consequence of the fact that
$\lambda _1a_1+\lambda _2a_2\leq \max\{a_1,a_2\}$ for any non-negative numbers
with $\lambda _1+\lambda _2=1$ (convex combination of the $a_i$'s).
Now taking limits as $h\searrow0$, we conclude that
\opteqn{
\mu (J)\leq 
\max \left\{-\cc_1 + \frac{\rho _2}{\rho _1} \kk, -\cc_2\right\}
=
-\cc \,,
}
as desired.

\mysubsection{Proof of Theorem~\protect{\ref{theo:fhc}}}
We first make some general remarks about perturbed systems.
Consider additive perturbations of the system~(\ref{eqn:gensys}) of the
following general form:
\be{eqn:perturbed}
\dot x \,=\, F(x,t) \,=\, f(t,x) + h(t,x)
\ee
where the vector field $h(t,x)$ is defined for $t\geq 0$ and $x\in C$, with values
in $\R^n$, is differentiable on $x$, and $h(t,x)$ and its Jacobian $H(t,x) =
\frac{\partial h}{\partial x}(t,x)$ are both continuous in $(t,x)$.
We have the following simple observation:

\begin{lemma}
\label{lem:perturb}
Assume that the system $\dot x=f(t,x)$ is infinitesimally contracting with
contraction rate $\cc$ with respect to a norm $\absv{\cdot }$.
Suppose that the Jacobian of the perturbation satisfies:
\be{eq:jacpert}
\norm{H(t,x)} \leq  c_h < c
\ee
for all $t\geq 0$ and all $x\in C$.
Then, the perturbed system~(\ref{eqn:perturbed})
is infinitesimally contracting with respect to the same norm.
\end{lemma}

\begin{proof}
The Jacobian of the new system is $\widetilde J(t,x) = J(t,x)+H(t,x)$, and:
\[
\mu (\widetilde J(x,t))
\,\leq \,
\mu \left(J\left(x,t\right)\right) +
\mu (\left(H\left(x,t\right)\right)
\,\leq \,
\widetilde c := -c + c_h
\]
by subadditivity of matrix measures and the fact that the norm always
upper-bounds the matrix measure (see for
instance~\cite{desoer-vidyasagar-siam}, page 31).
\end{proof}

Some comments regarding Lemma~\ref{lem:perturb} are as follows.
\emph{(i)}
Suppose that $h(t,x)$ does not depend on $x$.  Then~(\ref{eq:jacpert})
is trivially satisfied ($c_h=0$).
\emph{(ii)}
Suppose that
\opteqn{
H(t,x) \rightarrow  0
}
as $t\rightarrow \infty $, uniformly on $x\in C$.
Then the system $\dot x=F(t-t_0,x)$ is infinitesimally contracting.
That is, for any
two solutions $x(t)=\varphi(t,t_0,\xi )$ and $z(t)=\varphi(t,t_0,\zeta )$ 
of~(\ref{eqn:perturbed}) starting at time $t_0$, we have that:
\[
\abs{x(t)-z(t)} \;\leq \; e^{-\cc (t-t_0)} \abs{\xi -\zeta }, 
\quad\quad\forall\,t\ge t_0\geq 0\,.
\]
Indeed, by assumption we have that
$\beta (t) := \sup_{x\in C}\norm{H\left(x,t\right)} \rightarrow 0$ ,
so we can pick any $t_0>0$ so that $c_h=\beta (t_0)<c$.
\emph{(iii)}
Consider any two solutions $x(t)=\varphi(t,0,\xi )$ and
$z(t)=\varphi(t,0,\zeta )$ starting at time $t=0$.  Since
$x(t)=\varphi(t,t_0,x(t_0))$ and $z(t)=\varphi(t,t_0,z(t_0))$,
it follows that $x(t)-z(t)\rightarrow 0$ as $t\rightarrow 0$ (but not necessarily 
satisfying an estimate $\abs{x(t)-z(t)} \;\leq \; e^{-\cc t} \abs{\xi -\zeta }$).

\begin{lemma}
\label{lem:fh}
Assume that the system $\dot x=f(t,x)$ is infinitesimally contracting with
contraction rate $\cc$ with respect to a norm $\absv{\cdot }$.
Suppose that $h$ and its Jacobian $H$ are exponentially decreasing, in the
sense that, for some $k>0$:
\opteqn{
h(t,x)\,e^{kt} \quad \mbox{is bounded}
}
and
\opteqn{
\norm{H(t,x)e^{kt}} \leq  c_h < c\quad
\forall \,x \in C, \, \forall \, t\geq 0\,.
}
Then, there exist constants $\ell>0$ and $\kappa $ such that the following property
holds:
for any solution $x(t)=\varphi(t,0,\xi )$ of the system $\dot x=f(t,x)$,
and any solution $z(t)=\varphi(t,0,\zeta )$ of the system $\dot x=f(t,x)+h(t,x)$,
the estimate~(\ref{eqn:estimatefh}) is valid for all $t\geq 0$.
\end{lemma}

In the special case that $h$ is independent of $x$, this
proves Theorem~\ref{theo:fhc}.

\begin{proof}
Consider the following auxiliary system, with $p\in [0,1]$:
\beqn
\dot p &=& -kp\\
\dot x &=& F_p(t,p,x) \;=\; f(t,x) \,+\, p\, h(t,x)e^{kt}
\eeqn
viewed as a cascade.
The $p$-subsystem is infinitesimally contracting
with respect to the standard norm in $\R$.
The $x$-subsystem is infinitesimally contracting
when $p$ is viewed as a parameter.  Indeed, with:
\opteqn{
C(t,p,x) = \frac{\partial F_p}{\partial x}(t,p,x) = J(t,x)+ pH(t,x)e^{kt} \,,
}
we have that $\mu (C(t,p,x)) \leq  -c+c_h$, as earlier.
Moreover, the mixed Jacobian
\opteqn{
B(t,x,y) = \frac{\partial F_p}{\partial p}(t,p,x) = h(t,x)e^{kt}
}
is bounded, by assumption.
It follows from Theorem~\ref{theo:cascades} that the auxiliary system is also
infinitesimally contracting with some rate $\ell$, and the proof of that
result shows that this contraction can be established with respect to a norm
of the form: 
$\abs{(p,x)} = \rho _1 \abs{p}_{\norma} + \rho _2 \abs{x}_{\normb}$
for some $\rho _1>0$ and $\rho _2>0$, where $\abs{p}_{\norma}$ denotes the usual
norm in $\R$ and $\abs{x}_{\normb}$ denotes the original norm on $x$.

Consider now any solution $x(t)=\varphi(t,0,\xi )$ of the system $\dot x=f(t,x)$ and any
solution $z(t)=\varphi(t,0,\zeta )$ of the system $\dot x=f(t,x)+h(t,x)$.

Introduce $X(t) := (0,x(t))$ and $Z(t) := (e^{-kt},z(t))$.
It is clear that $X(t)$ and $Z(t)$ are the solutions of the
auxiliary system corresponding to initial conditions
$X(0)=(0,\xi )$ and $Z(0)=(1,\zeta )$ respectively.
Because the auxiliary system is infinitesimally contracting,
\opteqn{
\abs{X(t)-Z(t)}\leq e^{-\ell t}\abs{X(0)-Z(0)}
}
for all $t\geq 0$, where
\opteqn{
\abs{X(t)-Z(t)} = \rho _1 e^{-kt} + \rho _2 \abs{x(t)-z(t)}_{\normb}
}
and
\opteqn{
\abs{X(0)-Z(0)} = \rho _1  + \rho _2 \abs{\xi -\zeta }_{\normb}\,.
}
So
\opteqn{
\rho _2 \abs{x(t)-z(t)}_{\normb} \leq  
e^{-\ell t} \left(\rho _1  + \rho _2 \abs{\xi -\zeta }_{\normb}\right).
}
Dividing by $\rho _2$ and dropping the subscript for norms, we have
(\ref{eqn:estimatefh}) with $\kappa =\rho _1/\rho _2$.
\end{proof}

\section{Synchronization}
\label{sec:synch}

We remark here on the use of contraction theory to show synchronization of
coupled systems,
based on the introduction of ``virtual dynamics'' by Slotine
and collaborators (see for example~\cite{russo_dibernardo_ieee_cas2009}).
For simplicity of notation, we consider time-invariant dynamics,
but the same considerations apply to time-dependent vector fields.

Suppose that we have two diffusion-interconnected identical systems:
\beqn
\dot y&=& f(y) + \gamma (z) - \gamma (y)\\
\dot z&=& f(z) + \gamma (y) - \gamma (z)
\eeqn
where we think of $\gamma $ as a coupling law and assume that $\gamma $ is globally
Lipschitz. 
Typically, $\gamma $ is linear, so that $\gamma (z) - \gamma (y) = D(z-y)$ for a matrix $D$
(which is often a diagonal matrix).
For example, suppose that the systems are linear:
$\dot y= Ay + D(z-y)$ and $\dot z= Az + D(y-z)$.
Each system is individually (when $D=0$) asymptotically stable if and only
$A$ is a Hurtwitz matrix (all eigenvalues have negative real part).
Using a change of variables $(y,z) \mapsto  (y-z,y+z)$, we may bring this system to
a block-diagonal form with blocks
$A$$-$$2D$ and $A$,
and thus it is clear that the interconnected system is 
asymptotically stable if and only both $A$ and $A-2D$ are Hurwitz matrices.
Moreover, the same proof (the first block corresponds to $y-z$) shows that
for synchronization ($y(t)-z(t) \rightarrow 0$) it is enough that $A-2D$ be a
Hurwitz matrix.

For general, not necessarily linear systems, if the system $\dot x=f(x)$ is
infinitesimally contracting, then the decoupled systems (obtained when $\gamma =0$)
each satisfies that all solutions converge to each other.

More interestingly, a synchronization result can be established as follows.
Consider the following ``virtual system'':
\be{eq:virtual}
\dot x = f(x) - 2\gamma (x) + h(t)
\ee
(a different system results for each fixed input $h(\cdot )$)
and suppose that the vector field $f - 2h$ is infinitesimally contracting.
Take a particular solution $(y(t),z(t))$ of the coupled system.
Then, $y(t)$ and $z(t)$ are two solutions of~(\ref{eq:virtual}), when we pick
$h(t) = \gamma (y(t)) + \gamma (z(t))$.
It follows that
$\abs{z(t)-y(t)} \leq e^{-\cc t} \rightarrow  0$ for some $c>0$,
showing that the $y$ and $z$ subsystems synchronize.
Observe that this fact did not require the contractivity of $f$, but only
that of $f-2\gamma $.

Still for this solution $(y(t),z(t))$ of the coupled system,
we now define $h(t) = \gamma (z(t)) - \gamma (y(t))$.
Using the assumption that $\gamma $ is globally Lipschitz, we have that
$\abs{w(t)}\leq  M \abs{z(t)-y(t)} \leq M e^{-\cc t}$, for some constant $M$.
Now, if $f$ is contracting, we note that the equation
satisfied by $y$ is $\dot x=f(x)+h(t)$.
As $h(t)$ is exponentially convergent to zero, Theorem~\ref{theo:fhc}
implies that $y(t)-x(t)\rightarrow 0$ as $t\rightarrow \infty $ for every solution of the system
$\dot x=f(x)$.
Pick any one particular such solution $x_0(\cdot )$.
Then, $y(t)-x_0(t)\rightarrow 0$.
We may repeat this argument for an arbitrary $(y(t),z(t))$, always comparing
to the same $x_0(\cdot )$.
In summary, we have the following conclusion: if both $f$ and $f-2\eta $ are
infinitesimally contracting (not necessarily with respect to the same norm),
then all solutions of the coupled system converge to the diagonal solution
$(x_0(t),x_0(t))$.

The preceding considerations make the following question natural: when
does contractivity of $f$ (which is sufficient to provide a stability property
for the isolated  systems) already imply contractivity of $f-2\gamma $ (so that
synchronization to the uncoupled solutions occurs)?

We provide next a condition for the case when every Jacobian $D=D(x)$
of $\gamma (x)$ is a diagonal non-negative definite matrix.
The question is, then, for the Jacobians $A=J(x)$: when does
$\mu (A)\leq c$ imply that also $\mu (A-2D)\leq c$?

Recall that a norm on $\R^n$ is said to be monotonic or ``axis oriented''
if the following property holds for any two vectors in $\R^n$:
\opteqn{
\abs{y_i} \leq  \abs{x_i}   \Rightarrow  \abs{y} \leq \abs{x}\,.
}
The usual norms ($L^2$, $L^1$, $L^\infty $) are monotonic, as is any new norm
of the type $\abs{x}_P = \abs{Px}$ for a diagonal positive definite matrix $P$,
if $\abs{\cdot }$ is monotonic.

Theorems 2 and 3 of \cite{bauer} say that the following properties are
equivalent: (1) the norm is monotonic, (2) $\absv{x}$ depends only on the
absolute values of the components of $x$, and (3) the associated operator norm
satisfies that $\norm{E} = \max_j\{E_{jj}\}$ for any diagonal matrix $E$.
So $\norm{I-hD} = \max_j\{1-hD_{jj}\} = 1-hd_{ii}$ for some $i$,
which implies that $(1/h)(\norm{I-hD}-1) = -d_{ii}$ and thus
$\mu (D) = d_{ii}\leq 0$.
From subadditivity of matrix measures, we conclude that, for monotonic norms,
\opteqn{
\mu (A+D) \leq  \mu (A) \leq  c
}
and thus, for monotonic norms, we get contractivity of $f-2\gamma $ from
that of $f$.


\begin{thebibliography}{10}

\bibitem{angeli-sontag-fc}
D.~Angeli and E.~D. Sontag.
\newblock Forward completeness, unboundedness observability, and their
  {L}yapunov characterizations.
\newblock {\em Systems and Control Letters}, 38:209--217, 1999.

\bibitem{bauer}
F.L. Bauer, J.~Stoer, and C.~Witzgall.
\newblock Absolute and monotonic norms.
\newblock {\em Numerische Mathematik}, 3:257--264, 1961.

\bibitem{dahlquist}
G.~Dahlquist.
\newblock {\em Stability and error bounds in the numerical integration of
  ordinary differential equations}.
\newblock Trans. Roy. Inst. Techn. (Stockholm), 1959.

\bibitem{demidovich}
B.P. Demidovich.
\newblock Dissipativity of a nonlinear system of differential equations (in
  {R}ussian).
\newblock {\em Vestnik Moscow State University, Ser. Mat. Mekh.}, 6:19--27,
  1961.

\bibitem{desoer-vidyasagar-siam}
C.A. Desoer and M.~Vidyasagar.
\newblock {\em Feedback Synthesis: Input-Output Properties}.
\newblock SIAM, Philadelphia, 2009.

\bibitem{Jou_04_a}
J.~Jouffroy and J.~J.~E. Slotine.
\newblock Methodological remarks on contraction theory.
\newblock In {\em 42nd Conf. Decision and Control}, pages 2537--2543. IEEE
  Press, 2004.

\bibitem{Loh_Slo_00}
W.~Lohmiller and J.~J.~E. Slotine.
\newblock Nonlinear process control using contraction theory.
\newblock {\em AIChe Journal}, 46:588--596, 2000.

\bibitem{lozinskii}
S.~M. Lozinskii.
\newblock Error estimate for numerical integration of ordinary differential
  equations. {I}.
\newblock {\em Izv. Vtssh. Uchebn. Zaved. Mat.}, 5:222--222, 1959.

\bibitem{michelbook}
A.~N. Michel, D.~Liu, and L.~Hou.
\newblock {\em Stability of Dynamical Systems: Continuous, Discontinuous, and
  Discrete Systems}.
\newblock Springer-Verlag (New-York), 2007.

\bibitem{pavlov}
A.~Pavlov, A.~Pogromvsky, N.~van~de Wouv, and H.~Nijmeijer.
\newblock Convergent dynamics, a tribute to {B}oris {P}avlovich {D}emidovich.
\newblock {\em Systems and Control Letters}, 52:257--261, 2004.

\bibitem{pavlov_book}
A.~Pavlov, N.~{van de Wouw}, and H.~Nijmeijer.
\newblock {\em Uniform output regulation of nonlinear systems: a convergent
  dynamics approach}.
\newblock Springer-Verlag, Berlin, 2005.

\bibitem{russo_dibernardo_ieee_cas2009}
G.~Russo and M.~di~Bernardo.
\newblock Contraction theory and the master stability function: linking two
  approaches to study synchronization in complex networks.
\newblock {\em IEEE Transactions on Circuit and Systems II}, 56:177--181, 2009.

\bibitem{russo_dibernardo_sontag}
G.~Russo, M.~{di Bernardo}, and E.D. Sontag.
\newblock Global entrainment of transcriptional systems to periodic inputs.
\newblock 2009.
\newblock Submitted. Preprint in http://arxiv.org/abs/0907.0017.

\bibitem{soderlind}
G.~S\"oderlind.
\newblock The logarithmic norm: history and modern theory.
\newblock {\em BIT Numerical Mathematics}, 46:631--652, 2006.

\bibitem{mct}
E.~D. Sontag.
\newblock {\em Mathematical Control Theory. Deterministic Finite-Dimensional
  Systems}.
\newblock Springer-Verlag (New York), 1998.

\bibitem{arxiv_chaos_2009}
E.D. Sontag.
\newblock An observation regarding systems which converge to steady states for
  all constant inputs, yet become chaotic with periodic inputs.
\newblock Technical report, arxiv 0906.2166, 2009.

\bibitem{strom}
T.~Strom.
\newblock On logarithmic norms.
\newblock {\em SIAM J. Numer. Anal.}, 12:741--753, 1975.

\bibitem{Vid_93}
M.~Vidyasagar.
\newblock {\em Nonlinear systems analysis (2nd Ed.)}.
\newblock Pretice-Hall (Englewood Cliffs, NJ), 1993.

\bibitem{Wan_Slo_05}
W.~Wang and J.~J.~E. Slotine.
\newblock On partial contraction analysis for coupled nonlinear oscillators.
\newblock {\em Biological Cybernetics}, 92:38--53, 2005.

\bibitem{yoshizawa1}
T.~Yoshizawa.
\newblock {\em Stability theory by Liapunov's second method}.
\newblock The Mathematical Society of Japan, Tokyo, 1966.

\bibitem{yoshizawa2}
T.~Yoshizawa.
\newblock {\em Stability theory and the existence of periodic solutions and
  almost periodic solutions}.
\newblock Springer-Verlag, New York, 1975.

\end{thebibliography}



\edo

